\section{Analytic oscillation and flat seeds}
\label{sec:flat-balls}

This section provides two inputs for the later geometric argument.
Large oscillations of the Riesz pairing occur on a Carleson family of
cubes, and, outside another Carleson family, every cube has a
bilaterally flat descendant at a uniformly bounded depth.  We first prove the oscillation estimate by a Bessel inequality for
mean-zero Lipschitz tests. For the second assertion, two successive
tangents give a local alternative: absence of a flat ball forces a large
annular transform. Differences of local averages then turn that
alternative into a packing estimate.

Throughout the section $1\leq n<d$, and $\mu$ satisfies the
AD regularity and operator hypotheses of Theorem~\ref{thm:main}.
We use the global upper growth constant $G$ from
Lemma~\ref{lem:global-growth}.  Constants may depend on
$d,n,C_{\rm AD},C_{\rm R}$ and the lattice constants; additional
parameters are indicated when they occur.

\subsection{Packing the oscillation tests}

We use the Lipschitz-test and finite-depth Haar Riesz-system estimates
developed in \cite[Section~14]{NazarovTolsaVolberg2014} and include
the estimates with our normalization below. The mean-zero Lipschitz
test estimate and its oscillation-packing consequence also appear in
\cite[Section~5 and Appendix~B]{JayeNazarov2013}.

Write $\mathcal D(Q_0)=\{Q\in\mathcal D:Q\subset Q_0\}$.
We use the Carleson condition \eqref{eq:carleson-definition}, the
oscillation quantity \eqref{eq:oscillation-definition}, and the cell
mass bounds from Proposition~\ref{prop:dyadic-lattice}:
\begin{equation}\label{eq:packing-cell-masses}
 c_1\ell(Q)^n\le\mu(Q)\le C_1\ell(Q)^n.
\end{equation}
For bounded support, the diameter-truncated lattice has at most one
level with $\ell(Q)>D_E$. That level may be included in an exceptional
family at Carleson cost one. No estimate uses infinitely many scales
larger than the support diameter. Put $E^*=E\setminus N$, where $N$ is
the common null set in Proposition~\ref{prop:dyadic-lattice}. Its
partitions and nesting hold exactly on $E^*$. Also $E^*$ is dense in
$E$, since every ball centered on $E$ has positive measure.

\begin{lemma}[Bessel estimate for dyadic tests]\label{lem:dyadic-test-bessel}
Fix $J>1$. For each cube in a finite family $\mathcal A\subset\mathcal D$,
choose one test $\varphi_Q$ from the class defining $W_J(Q)$. There is
$B_J<\infty$, independent of $\mathcal A$ and the choices, such that
\begin{equation}\label{eq:dyadic-test-bessel}
 \sum_{Q\in\mathcal A}
 \frac{\left|\int\varphi_Q u\,d\mu\right|^2}{\ell(Q)^n}
 \le B_J\int |u|^2\,d\mu\qquad(u\in L^2(\mu)).
\end{equation}
For vector-valued $u$ the same constant works with
$\int\varphi_Q u$ replaced by
$\int\varphi_Q\langle e_Q,u\rangle$, where each $|e_Q|\le1$ may be chosen
independently.
\end{lemma}

\begin{proof}
Put $r_Q=\ell(Q)$, $w_Q=r_Q^{n/2}$ and $f_Q=\varphi_Q/w_Q$.
The support and Lipschitz conditions give $|\varphi_Q|\le2J$: compare a
point of its support with a point on the sphere of radius $Jr_Q$ where
the function vanishes. The tests are in $L^1\cap L^2$ by compact support.
If $r_P\le r_Q$, cancellation in $f_P$ gives
\begin{align}
 |\langle f_Q,f_P\rangle|
 &=\left|\int(f_Q(x)-f_Q(z_P))f_P(x)\,d\mu(x)\right|\notag\\
 &\le H_J\frac{r_P}{r_Q}\frac{w_P}{w_Q},
 \qquad H_J=GJ^{n+1}(2J+1)^2.
 \label{eq:gram-cancellation}
\end{align}
Indeed, on the support of $f_P$ the difference is at most
$Jr_P/(r_Qw_Q)$, and
\[
 \int|f_P|\,d\mu\le \frac{2JG(Jr_P)^n}{w_P}.
\]
The symmetric estimate follows by reversing $P,Q$.

A nonzero inner product requires the two support balls to meet. Fix
$Q$ and a finer level $r_P=2^{-k}r_Q$. Every interacting cube $P$ lies
in $B(z_Q,(2J+C_0)r_Q)$, up to its null boundary. Same-level cells are
pairwise disjoint, so \eqref{eq:packing-cell-masses} and global upper
growth imply
\begin{equation}\label{eq:gram-fine-mass}
 \sum_{\substack{P\text{ at this level}\\
                  \langle f_Q,f_P\rangle\ne0}}r_P^n
 \le H'_J r_Q^n,
 \qquad H'_J=c_1^{-1}G(2J+C_0)^n.
\end{equation}
At a coarser level $r_P=2^kr_Q$, the same argument in
$B(z_Q,(2J+C_0)r_P)$ bounds the number of interacting cubes by $H'_J$.
These are estimates in the measure dimension $n$.

Multiply \eqref{eq:gram-cancellation} by $w_P$. Summing first over a
finer level using \eqref{eq:gram-fine-mass}, and then over a coarser level
using its counting bound, gives respectively
\[
 \sum_{r_P=2^{-k}r_Q}|\langle f_Q,f_P\rangle|w_P
 \le H_JH'_J2^{-k}w_Q,
 \qquad
 \sum_{r_P=2^kr_Q}|\langle f_Q,f_P\rangle|w_P
 \le H_JH'_J2^{-k}w_Q.
\]
Zero inner products may be discarded. Summing the geometric series
therefore yields
\[
 \sum_{P\in\mathcal A}|\langle f_Q,f_P\rangle|w_P
 \le4H_JH'_Jw_Q.
\]
For completeness, a symmetric matrix $(a_{QP})$ of nonnegative entries
with $\sum_P a_{QP}w_P\le Bw_Q$ satisfies
\[
 \sum_{Q,P}a_{QP}|t_Qt_P|\le B\sum_Q|t_Q|^2.
\]
This follows by applying
$2|t_Qt_P|\le |t_Q|^2w_P/w_Q+|t_P|^2w_Q/w_P$ and using symmetry.
Apply it to $a_{QP}=|\langle f_Q,f_P\rangle|$. The resulting bound for
$\|\sum_Qt_Qf_Q\|_2^2$ gives \eqref{eq:dyadic-test-bessel} by Hilbert-space
duality, with $B_J=4H_JH'_J$.

For the vector assertion, write $v_Q=\int f_Qu\,d\mu$. Then
$|\langle e_Q,v_Q\rangle|^2\le|v_Q|^2$. Expand the latter in an
orthonormal basis, apply the scalar estimate to each coordinate of $u$,
and sum by Parseval's identity. The constant is unchanged.
\end{proof}

\begin{proposition}[Oscillation packing]\label{prop:oscillation-packing}
For every $J>1$ and $v>0$, the family
\[
 \mathcal O(J,v)=\{Q\in\mathcal D:W_J(Q)>v\ell(Q)^n\}
\]
is Carleson. Its constant depends only on the fixed data, $J$, and $v$.
\end{proposition}

\begin{proof}
Fix $Q_0$, put $L=\ell(Q_0)$, and fix the single outer ball
\[
 S=B(z_{Q_0},(C_0+2)L).
\]
Let $\mathcal A\subset\mathcal O(J,v)\cap\mathcal D(Q_0)$ be finite.
For each selected cube write $r_Q=\ell(Q)$ and choose a witnessing test
$\varphi_Q$ and a unit-bounded direction $e_Q$, so that
\[
 |\langle e_Q,\mathcal R_\mu(\varphi_Q)\rangle|>vr_Q^n.
\]
For $r_Q\le L/(2J)$ the support of the test lies in $S$, and its distance
from $S^c$ is at least $L$. With $\nu=\mu|S$, cancellation gives
\begin{align*}
 \mathcal R_\mu(\varphi_Q)-\mathcal R_\nu(\varphi_Q)
 =\int\varphi_Q(x)\int_{S^c}
 [K_n(x-y)-K_n(z_Q-y)]\,d\mu(y)\,d\mu(x).
\end{align*}
This integral is absolutely convergent. The kernel derivative bound and
upper growth on dyadic annuli give
\[
 \int_{|y-z_Q|\ge L}|y-z_Q|^{-n-1}\,d\mu(y)
 \le C(n)G/L.
\]
Since $|x-z_Q|\le Jr_Q$ on the support and
$\int|\varphi_Q|\,d\mu\le2JG(Jr_Q)^n$, it follows that
\begin{equation}\label{eq:oscillation-outer-error}
 |\mathcal R_\mu(\varphi_Q)-\mathcal R_\nu(\varphi_Q)|
 \le T_J(r_Q/L)r_Q^n,
 \qquad T_J=C(n)G^2J^{n+2}.
\end{equation}
Choose an integer $N\ge1$, depending only on $J,v$ and the fixed data,
with
\[
 2^{-N}\le\min\{(2J)^{-1},\ v/[2(T_J+1)]\}.
\]
If $Q$ lies at least $N$ generations below $Q_0$, its restricted pairing
therefore has magnitude at least $(v/2)r_Q^n$ in the direction $e_Q$.

The same hard-truncated function $R_{\mu,\varepsilon}1_S$ can be used for
all these fine cubes. For each such cube,
\[
 \int\varphi_Q(x)R_{\mu,\varepsilon}1_S(x)\,d\mu(x)
 =\frac12\iint_{S\times S,\ |x-y|>\varepsilon}
  (\varphi_Q(x)-\varphi_Q(y))K_n(x-y)\,d\mu(x)d\mu(y).
\]
As $\varepsilon\downarrow0$ this converges to
$\mathcal R_\nu(\varphi_Q)$. Antisymmetry gives the displayed identity;
Lipschitz cancellation and upper growth make its near-diagonal majorant
integrable. Also $\int\varphi_Q\,d\nu=0$ because the support lies in $S$.
Lemma~\ref{lem:dyadic-test-bessel} and the hard-truncation bound imply
\[
 \sum_{Q\in\mathcal A}
 \frac{\left|\int\varphi_Q\langle e_Q,
                R_{\mu,\varepsilon}1_S\rangle\,d\mu\right|^2}{r_Q^n}
 \le B_J C_{\rm R}^2\mu(S).
\]
Restrict the sum to the fine cubes and pass to the limit in this finite
sum. We obtain
\[
 \sum_{\substack{Q\in\mathcal A\\r_Q\le2^{-N}L}}r_Q^n
 \le4v^{-2}B_JC_{\rm R}^2\mu(S).
\]
Now $\mu(S)\le G(C_0+2)^nL^n\le
G(C_0+2)^nc_1^{-1}\mu(Q_0)$, so
\eqref{eq:packing-cell-masses} bounds the total mass of the fine cubes
by a fixed multiple of $\mu(Q_0)$. At each of the fewer than $N$ remaining
levels, disjointness and containment in $Q_0$ bound the selected masses
by $\mu(Q_0)$. This proves the desired estimate for every finite
$\mathcal A$, hence for the entire family.
\end{proof}

\subsection{Flat descendants outside a Carleson family}

The second packing estimate provides a small flat region inside every
cube outside a Carleson family. The allowed depth depends on the desired
aperture and flatness, but is uniform over all cubes.

\begin{proposition}[Flat seeds outside a Carleson family]\label{prop:flat-seeds}
For every $A_{\rm F}\ge1$ and $\eta>0$ there are an integer $N_{\rm F}\ge1$
and a Carleson family $\mathcal F$ such that every
$Q\in\mathcal D\setminus\mathcal F$ has a proper descendant $P\subset Q$
with
\begin{equation}\label{eq:flat-seed-conclusion}
 1\le\operatorname{depth}_Q(P)\le N_{\rm F},
 \qquad b_\mu(z_P,A_{\rm F}\ell(P))<\eta.
\end{equation}
One may also require
$B(z_P,A_{\rm F}\ell(P))\subset B(z_Q,c_0\ell(Q)/4)$.
The constants are uniform in $Q$ and the top cube used in the Carleson
estimate; they depend only on the fixed data, $A_{\rm F}$ and $\eta$.
\end{proposition}

We prove this proposition by detecting a large annular transform whenever
such a descendant is absent. The next lemma supplies the geometric
alternative independently of the operator bound; the subsequent
comparison of averages makes its annular branch available to a Bessel
estimate.

\subsubsection{The local annular alternative}

The argument follows the flat-ball/large-annulus method of Nazarov,
Tolsa and Volberg; see \cite[Section~10]{NazarovTolsaVolberg2014}.
We give the version needed here in arbitrary codimension, with the
stated smooth annuli and admissible-radius conventions.

Fix a smooth radial function $\chi:\R^d\to[0,1]$ which equals one
on $\overline B(0,1)$ and vanishes outside $B(0,2)$.  For $r>0$ put
$\chi_{a,r}(y)=\chi((y-a)/r)$, and, for $0<r\leq R$, define
\begin{equation}\label{eq:smooth-annular-transform}
 \mathcal S_\mu(a;r,R)
   =\int \bigl(\chi_{a,R}(y)-\chi_{a,r}(y)\bigr)
                 K_n(a-y)\,d\mu(y).
\end{equation}
This is an absolutely convergent vector integral: its integrand is
bounded, compactly supported, and zero near $a$.  The same definition
will be used for any measure with locally finite mass.

\begin{lemma}[A flat ball or a large annulus]\label{lem:annular-alternative}
Fix $0<\kappa<1/8$, $\varepsilon>0$, and $T>0$.  There exists
$\rho\in(0,\kappa)$, depending only on
$d,n,C_{\rm AD},G,\kappa,\varepsilon,T$, with the following property.
For every $a_0\in E$ and $0<s\leq D_E$, at least one of the
following assertions holds:
\begin{enumerate}
\item There are $a\in E\cap B(a_0,\kappa s)$ and
      $\rho s\leq t<\kappa s$ such that
      $b_\mu(a,t)<\varepsilon$.
\item There are $a\in E\cap B(a_0,\kappa s)$ and
      $\rho s\leq r\leq R<\kappa s$ such that
      $|\mathcal S_\mu(a;r,R)|>T$.
\end{enumerate}
The number $\rho$ is chosen before $\mu$, $a_0$, and $s$.
The operator bound is not needed in this lemma.
\end{lemma}

\begin{proof}
We first prove a qualitative assertion for measures with global
upper and lower growth.  We then localize it by taking a tangent and
obtain the uniform radius floor by compactness.

\paragraph{Support convergence and annular bounds.}
We record the elementary convergence facts needed at both stages.
Suppose $\sigma_j$ converge weakly on compact sets to $\sigma$,
have a common upper growth bound, and satisfy
\begin{equation}\label{eq:capped-lower-growth-flat}
 \sigma_j(B(x,t))\geq c t^n
 \qquad(x\in\supp\sigma_j,\quad 0<t\leq h),
\end{equation}
where $c,h>0$ are fixed.  On every bounded region their supports
converge in both directions, with an arbitrarily small enlargement
of that region.  Indeed, if bounded support points $x_j$ stayed a
positive distance from $\supp\sigma$, a subsequence would converge
to such a point $x$.  A continuous cutoff near $x$, together with
\eqref{eq:capped-lower-growth-flat} at a sufficiently small fixed
radius, would have uniformly positive $\sigma_j$ integrals and
zero $\sigma$ integral.  Conversely, a ball about a point of
$\supp\sigma$ contains a nonnegative compactly supported
continuous test with positive $\sigma$ integral.  Its $\sigma_j$
integral is eventually positive.  Finite coverings make both
assertions uniform on bounded sets.  The limit also has lower growth:
for $x\in\supp\sigma$ and $0<t\leq h$, choose $x_j\to x$
and a continuous cutoff equal to one on $B(x,t/2)$ and supported
in $B(x,t)$.  Eventually $B(x_j,t/4)\subset B(x,t/2)$, so passing
the cutoff integrals to the limit gives
$\sigma(B(x,t))\geq c(t/4)^n$.

In particular, if $a_j\in\supp\sigma_j$ tend to
$a\in\supp\sigma$ and the limit is sufficiently flat on
$B(a,2t)$, then the sources are flat on $B(a_j,t)$.
For example, for $0<\eta\leq1$,
\begin{equation}\label{eq:nonflat-limit-flat}
 \left[ b_{\sigma_j}(x,t)\geq\eta
 \text{ for all nearby support centers and all large }j\right]
 \quad\Longrightarrow\quad
 b_\sigma(a,2t)\geq\eta/64.
\end{equation}
To check the numerical slack, if the last inequality failed, choose
an $n$-plane whose sum of deviations on $B(a,2t)$ is
less than $\eta t/32$.  Approximate $a$ by $a_j$ and both supports
within $\eta t/32$.  Points of either set in $B(a_j,t)$ and all
the points used to approximate them then lie in $B(a,2t)$.
Each of the two deviations for the sources is less than
$\eta t/8$, contradicting $b_{\sigma_j}(a_j,t)\geq\eta$.
This reasoning uses only the lower growth at small fixed radii;
the tested radius $t$ need not be below $h$.

For fixed $r,R>0$, the kernel in
\eqref{eq:smooth-annular-transform} is continuous with compact
support.  If its center varies from $a_j$ to $a$, the kernels
converge uniformly and have a common bounded support.  The upper
growth bound therefore gives
\begin{equation}\label{eq:annular-center-convergence}
 \mathcal S_{\sigma_j}(a_j;r,R)
       \longrightarrow \mathcal S_\sigma(a;r,R).
\end{equation}
Here and below support centers in a limit can be approximated by
source support centers, as just proved.

For $u>0$ define the normalized dilation
\[
 \sigma_{a,u}=u^{-n}
       \bigl(y\mapsto (y-a)/u\bigr)_\#\sigma.
\]
A change of variables gives the exact identities
\begin{equation}\label{eq:annular-flat-scaling}
 \begin{split}
 \mathcal S_{\sigma_{a,u}}(b;r,R)
      &=\mathcal S_\sigma(a+ub;ur,uR),\\
 b_{\sigma_{a,u}}(b,t)&=b_\sigma(a+ub,ut).
 \end{split}
\end{equation}
Thus a bound $|\mathcal S_\sigma(x;r,R)|\leq T$ at all support
centers and all positive radii survives dilations and compact
limits.  A global lower bound $b_\sigma\geq\eta$ survives a
compact limit with the loss $\eta\mapsto\eta/64$.
The compactness and nonvanishing of normalized dilations follow
from Lemma~\ref{lem:weak-compactness}: upper growth controls the
mass on each compact set, and lower growth preserves positive mass
in every ball about the origin.

\paragraph{Two tangents reduce the containing dimension.}
Suppose, towards a contradiction, that a nonzero measure $\sigma$
satisfies global bounds
\[
 c t^n\leq\sigma(B(x,t)),\qquad
 \sigma(B(y,t))\leq G_1t^n
 \quad(x\in\supp\sigma,\ y\in\R^d,\ t>0),
\]
that $|\mathcal S_\sigma(x;r,R)|\leq T$ at every support
center and all $0<r\leq R$, and that
$b_\sigma(x,t)\geq\eta>0$ everywhere.  Decreasing $\eta$ if
necessary, assume $\eta\leq1$.  Let $P$ be a linear subspace
containing the support, initially $P=\R^d$.

The support cannot fill $P$.  If it did, upper and lower growth
would force $\dim P=n$.  To see this, write $k=\dim P$ and
cover its unit ball by $O(u^{-k})$ balls of radius $u$.
For $k<n$, upper growth would make its mass tend to zero.
For $k>n$, a family of $\gtrsim u^{-k}$ disjoint balls of
radius comparable to $u$, centered in the unit ball of $P$,
would have total mass $\gtrsim u^{n-k}$, contradicting the
upper bound on a fixed larger ball.  When $k=n$, the support
is an $n$-plane and all its bilateral beta numbers are zero.
Each possibility contradicts the assumptions.

Choose $c_1\in P\setminus\supp\sigma$ and a nearest support
point $a$ to $c_1$.  Such a point exists because the support is
nonempty and closed in finite-dimensional Euclidean space.
Set $e=a-c_1\in P\setminus\{0\}$.  The nearest-point inequality
implies
\[
 2\langle e,y-a\rangle+|y-a|^2\geq0
       \qquad(y\in\supp\sigma).
\]
After dilation at $a$ by scales $u_j\downarrow0$ this becomes
$2\langle e,x\rangle+u_j|x|^2\geq0$.
Pass to a nonzero limit $\tau$.  The support-convergence argument
shows that
\[
 \supp\tau\subset P\cap\{x:\langle e,x\rangle\geq0\},
 \qquad 0\in\supp\tau.
\]
This first tangent retains global upper and lower growth,
the annular bound $T$, and positive global nonflatness.

We next use the sign of the normal kernel to take a second
tangent.  Define
\[
 p_e(x)=\frac{\langle e,x\rangle_+}{|x|^{n+1}}
 \quad(x\ne0),\qquad p_e(0)=0,
 \qquad t_+=\max\{t,0\}.
\]
For $2r\leq R_0$, the weight
$\chi_{0,R_0}-\chi_{0,r}$ is nonnegative everywhere and equals
one when $2r\leq|x|\leq R_0$.  In fact, inside the $R_0$-ball
the outer cutoff is one, and outside it the inner cutoff is zero.
This argument requires no monotonicity of the cutoff.
On the support of $\tau$,
\[
 \int(\chi_{0,R_0}-\chi_{0,r})p_e\,d\tau
       =-\langle e,\mathcal S_\tau(0;r,R_0)\rangle
       \leq |e|T.
\]
Given $v>0$, take $r=\min(v,R_0)/4$ to obtain
\[
 \int_{v<|x|<R_0}p_e(x)\,d\tau(x)\leq |e|T.
\]
Monotone convergence, using $p_e(0)=0$, proves
$p_e\in L^1(\tau|_{B(0,R_0)})$.  This establishes an actual
integrability statement before taking the next limit.

For any fixed $R>0$, exact scaling gives
\begin{align}\label{eq:vanishing-normal-moment}
 \int_{B(0,R)}\langle e,x\rangle_+\,d\tau_{0,u}(x)
  &=u^{-n-1}\int_{B(0,uR)}\langle e,y\rangle_+\,d\tau(y)\notag\\
  &\leq R^{n+1}\int_{B(0,uR)}p_e(y)\,d\tau(y)
       \longrightarrow0 \qquad(u\downarrow0).
\end{align}
The convergence follows from absolute continuity of the integral:
upper growth and $n>0$ give $\tau(B(0,uR))\to0$.
Take a nonzero compact limit $\nu$ of these second dilations.
Testing with nonnegative compact continuous cutoffs times
$\langle e,x\rangle_+$ shows that this continuous nonnegative
function vanishes throughout $\supp\nu$.
The negative normal part already vanishes on every dilated
support.  Hence
\[
 \supp\nu\subset Q:=P\cap e^\perp,
 \qquad \dim Q<\dim P.
\]
The last inequality uses $e\in P\setminus\{0\}$:
$e$ itself is not in $e^\perp$.  Both tangents retain support
in $P$ because their centers lie in $P$.

The new measure is nonzero, satisfies global upper and lower
growth, has annular bound $T$, and has global beta lower bound
at least $\eta/4096$.  For definiteness, the compactness
estimates may multiply the upper constant by $2^n$ and divide
the lower constant by $4^n$ at each tangent; their precise
values are immaterial.  Strong induction on $\dim P$, allowing
these constants and the positive beta threshold to vary, gives
a contradiction.  We have proved that a globally regular,
globally annular-bounded measure has a flat ball at every
prescribed positive tolerance.

\paragraph{Localization and a uniform radius floor.}
Suppose a measure has global upper growth, lower growth only
for $0<t\leq h$, and annular bound $T$ for support centers in
an open set $U$ with $U\cap\supp\sigma\ne\varnothing$,
and outer radii below a positive cap.
If there were no $\eta$-flat ball in $U$ below another positive
cap, take a tangent at any point of $U\cap\supp\sigma$.
Every fixed normalized radius eventually lies below the physical
caps, and bounded normalized centers dilate into $U$.
Equations~\eqref{eq:annular-center-convergence}--
\eqref{eq:annular-flat-scaling} and
\eqref{eq:nonflat-limit-flat} therefore make the annular bound
and a positive beta lower bound global on the tangent.
The lower growth also becomes global, since its normalized cap
is $h/u_j\to\infty$.  This contradicts the preceding paragraph.
Thus the local qualitative alternative holds with no radius floor.

Normalize the statement of the lemma by $a_0=0$, $s=1$.
The resulting measures have common global upper growth, origin
in their supports, and lower growth at all support centers for
$0<t\leq1$.  If no common $\rho$ worked, choose counterexamples
$\mu_j$ with floors $\rho_j\downarrow0$.
The uniform mass bounds on compact sets give a locally weakly
convergent subsequence, by the compactness argument in
Lemma~\ref{lem:weak-compactness}.  The fixed lower bound at zero
keeps its limit $\nu$ nonzero and puts $0\in\supp\nu$.
The support-convergence argument above transfers the lower bound
at radii below a fixed positive cap to $\nu$.  Every fixed annulus with center in $B(0,\kappa)$ and
$0<r\leq R<\kappa$ eventually lies above the floor $\rho_j$.
Its transform on the limit is therefore bounded by $T$.
For every fixed $0<t<\kappa$, the failure of flat balls at
radius $t/2$ and \eqref{eq:nonflat-limit-flat} give
\[
 b_\nu(a,t)\geq\min(\varepsilon,1)/64
 \qquad(a\in\supp\nu\cap B(0,\kappa)).
\]
The local qualitative alternative contradicts these two
properties.  This proves the existence of $\rho$ before the
measure is chosen.  Finally,
\eqref{eq:annular-flat-scaling} transports the conclusion back
to $a_0,s$.  Only the original lower growth below $s\leq D_E$
has been used.
\end{proof}

\subsubsection{From large annuli to flat descendants}

The passage from a large annulus to a difference of local averages,
and then to a packing estimate, follows the mechanism in
\cite[Section~15]{NazarovTolsaVolberg2014}.

The next comparison makes the large-annulus branch of
Lemma~\ref{lem:annular-alternative} usable in an $L^2$ packing argument.
For a positive finite-measure set $A$, write
$\langle u\rangle_A=\mu(A)^{-1}\int_Au\,d\mu$.

\begin{lemma}[Annuli and differences of averages]\label{lem:annular-average-gap}
Let $0<r\le R$, let $S$ be a bounded measurable set containing
$B(a,2R)$, and let $A\subset B$ be measurable sets of positive finite
measure satisfying
\[
 A\subset B(a,r/2),\qquad B\subset B(a,R/2),\qquad
 r^n\le M\mu(A),\qquad R^n\le M\mu(B).
\]
There is a constant
\[
 E_0=3C_{\rm R}(G2^nM)^{1/2}+3C(n)G,
\]
independent of $r,R,R/r,S$ and $\varepsilon$, such that, for every
$0<\varepsilon<r/2$,
\begin{equation}\label{eq:annular-average-error}
 \left|\mathcal S_\mu(a;r,R)
 -\bigl(\langle R_{\mu,\varepsilon}1_S\rangle_A
       -\langle R_{\mu,\varepsilon}1_S\rangle_B\bigr)\right|
 \le E_0.
\end{equation}
In particular, one direction chosen from $\mathcal S_\mu(a;r,R)$
witnesses an average gap at least $|\mathcal S_\mu(a;r,R)|-E_0$
for every such $\varepsilon$.
\end{lemma}

\begin{proof}
Decompose the one fixed input as
\[
 1_S=f+g+h,\qquad
 f=\chi_{a,r},\quad g=\chi_{a,R}-\chi_{a,r},\quad
 h=1_S-\chi_{a,R}.
\]
Here $|g|,|h|\le1$, $g$ vanishes on $B(a,r)$, and $h$ vanishes on
$B(a,R)$. All three functions belong to $L^2(\mu)$. Upper growth gives
\[
 \|f\|_2^2\le G(2r)^n,
 \qquad \|g\|_2^2\le G(2R)^n.
\]
Cauchy--Schwarz and the operator bound consequently bound each of
\[
 |\langle R_{\mu,\varepsilon}f\rangle_A|,
 \quad |\langle R_{\mu,\varepsilon}f\rangle_B|,
 \quad |\langle R_{\mu,\varepsilon}g\rangle_B|
\]
by $C_{\rm R}(G2^nM)^{1/2}$; use $r\le R$ for the middle term.

If $u\in L^2(\mu)$, $|u|\le1$, and $u$ vanishes on $B(a,t)$, then for $|x-a|\le t/2$
and $0<\varepsilon<t/2$ there is no truncation boundary on its support.
The kernel derivative estimate therefore gives
\[
 |R_{\mu,\varepsilon}u(x)-R_{\mu,\varepsilon}u(a)|
 \le C(n)|x-a|\int_{|y-a|\ge t}|y-a|^{-n-1}\,d\mu(y)
 \le C(n)G.
\]
Apply this once to $g$ on $A$, and twice to $h$ on $A$ and $B$.
Since $R_{\mu,\varepsilon}g(a)=\mathcal S_\mu(a;r,R)$, expanding the
two averages of $R_{\mu,\varepsilon}(f+g+h)$ proves
\eqref{eq:annular-average-error}. Finally choose a unit-bounded $e$ with
$\langle e,\mathcal S_\mu(a;r,R)\rangle=|\mathcal S_\mu(a;r,R)|$.
This choice precedes $\varepsilon$; projecting
\eqref{eq:annular-average-error} onto $e$ proves the last assertion.
\end{proof}



\begin{proof}[Proof of Proposition~\ref{prop:flat-seeds}]
Put $L_0=\max\{1,2C_0\}$ and $M=(8L_0)^n/c_1$.
Fix the comparison constant $E_0$ in
Lemma~\ref{lem:annular-average-gap}, with this value of $M$.
Apply Lemma~\ref{lem:annular-alternative} with
\[
 \kappa=c_0/32,\qquad
 \epsilon_{\rm F}=\frac{\eta A_{\rm F}}{32(A_{\rm F}+C_0)},
 \qquad T=E_0+2.
\]
Let $\rho>0$ be its uniform relative radius. These choices, including
$\rho$, are made before the cube $Q$ is chosen. For now assume
$\ell(Q)\le D_E$ and apply that lemma at $a_0=z_Q$, $s=\ell(Q)$.

\smallskip\noindent
\emph{The flat branch produces a seed.}
Suppose the lemma gives $a\in E\cap B(z_Q,\kappa\ell(Q))$ and
$\rho\ell(Q)\le t<\kappa\ell(Q)$ with $b_\mu(a,t)<\epsilon_{\rm F}$.
Choose $b\in E^*\cap B(a,t/8)$, and choose the dyadic scale of $P$ so that
\[
 \frac{t}{16(A_{\rm F}+C_0)}<\ell(P)
 \le\frac{t}{8(A_{\rm F}+C_0)}.
\]
Let $P$ be the cell at that scale containing $b$.
Both $b\in Q$ and $\ell(P)<\ell(Q)$ follow from the core containment
and $t<\kappa\ell(Q)$, so exact nesting on $E^*$ gives $P\subset Q$.
Also $|z_P-b|\le C_0\ell(P)$, whence
\[
 B(z_P,A_{\rm F}\ell(P))\subset B(a,t)
 \cap B(z_Q,c_0\ell(Q)/4).
\]
Choose an affine plane for which the sum of the two deviations in
$B(a,t)$ is less than $\epsilon_{\rm F}t$. Restricting both suprema to
the smaller ball gives
\[
 b_\mu(z_P,A_{\rm F}\ell(P))
 <\frac{\epsilon_{\rm F}t}{A_{\rm F}\ell(P)}
 <\frac{16\epsilon_{\rm F}(A_{\rm F}+C_0)}{A_{\rm F}}
 =\eta/2.
\]
The lower bound for $t$ bounds the depth of $P$ uniformly. For example,
any integer $N_{\rm F}>\log_2(16(A_{\rm F}+C_0)/\rho)$ suffices.

\smallskip\noindent
\emph{Cubes without a flat seed have a large annulus.}
Let $\mathcal F_1$ be the cubes with $\ell(Q)\le D_E$ for which no seed
as in \eqref{eq:flat-seed-conclusion}, with the additional ball
containment, exists at depth at most $N_{\rm F}$. They must lie in the
annular branch. Thus each has radii
\[
 \rho\ell(Q)\le r_Q\le R_Q<\kappa\ell(Q)
\]
and a center $a_Q$ in the same core with
$|\mathcal S_\mu(a_Q;r_Q,R_Q)|>E_0+2$.
For these fixed positive radii the annular transform is continuous in
its center. Indeed the annular kernel is zero near its center, is
continuous there after extension by zero, and has bounded compact
support locally uniformly in the center; dominated convergence applies
using local finiteness. Density of $E^*$ therefore permits replacing
$a_Q$ by $a'_Q\in E^*$ so close that
\[
 |a'_Q-z_Q|<2\kappa\ell(Q),\qquad
 |\mathcal S_\mu(a'_Q;r_Q,R_Q)|>E_0+1.
\]
This perturbation uses the strict core and norm margins, and does not
change either radius.

Choose dyadic scales $s_r,s_R$ with
\[
 \frac{r_Q}{8L_0}<s_r\le\frac{r_Q}{4L_0},\qquad
 \frac{R_Q}{8L_0}<s_R\le\frac{R_Q}{4L_0},
\]
using the same monotone rounding rule, so $s_r\le s_R$.
Let $A_Q$ and $B_Q$ be the cells at those scales containing $a'_Q$.
They are descendants of $Q$ with $A_Q\subset B_Q\subset Q$.
Since a cell at scale $s$ has diameter at most $L_0s$,
\[
 A_Q\subset B(a'_Q,r_Q/2),\qquad
 B_Q\subset B(a'_Q,R_Q/2),
 \quad r_Q^n\le M\mu(A_Q),\quad R_Q^n\le M\mu(B_Q).
\]
Fix an integer $I\ge1$ larger than $\log_2(8L_0/\rho)$.
The depth of $A_Q$ below $Q$ is at most $I$, and
\begin{equation}\label{eq:haar-inner-parent-mass}
 \mu(Q)\le H_I\mu(A_Q),\qquad
 H_I=(C_1/c_1)2^{In}.
\end{equation}
Neither $M$ nor $E_0$ depends on $\rho$; the depth cost $H_I$ is allowed
to depend on the now fixed $\rho$.

\smallskip\noindent
\emph{Haar Bessel gives the Carleson estimate.}
Set
\[
 h_Q=\mu(A_Q)^{1/2}
       \left(\frac{1_{A_Q}}{\mu(A_Q)}-
             \frac{1_{B_Q}}{\mu(B_Q)}\right).
\]
Direct integration gives $\int h_Q\,d\mu=0$ and
$\|h_Q\|_2^2=1-\mu(A_Q)/\mu(B_Q)\le1$.
This function is supported in $Q$ and is constant almost everywhere on
every cell at least $I$ generations below $Q$.
Partition a finite collection of these functions according to the
level of $Q$ modulo $I$. Within one class, distinct equal-level parents
are disjoint. For unequal levels, the gap is at least $I$: the coarser
function is constant on the finer parent, or their supports are
disjoint. The finer function's mean zero gives orthogonality. Bessel's
inequality in each class and summation give
\begin{equation}\label{eq:bounded-depth-haar-bessel}
 \sum_Q\left|\int h_Q\langle e_Q,u\rangle\,d\mu\right|^2
 \le I\|u\|_2^2
\end{equation}
for vector $u\in L^2(\mu)$ and independently chosen $|e_Q|\le1$.
The vector extension follows by the same Parseval argument as in
Lemma~\ref{lem:dyadic-test-bessel}.

Fix a top cube $Q_0$ and a finite
$\mathcal A\subset\mathcal F_1\cap\mathcal D(Q_0)$. Use, for every
selected cube, the same set
$S=B(z_{Q_0},(C_0+2)\ell(Q_0))$.
The annular center and radius bounds ensure
$B(a'_Q,2R_Q)\subset S$.
Choose $e_Q$ from the annular vector as in
Lemma~\ref{lem:annular-average-gap}; then choose a single
\[
 0<\varepsilon<\tfrac12\min_{Q\in\mathcal A}r_Q
\]
when the family is nonempty. The average gap is greater than one for
every selected $Q$ at this same truncation. Consequently
\[
 \mu(A_Q)
 <\left|\int h_Q
       \langle e_Q,R_{\mu,\varepsilon}1_S\rangle\,d\mu\right|^2.
\]
By \eqref{eq:haar-inner-parent-mass},
\eqref{eq:bounded-depth-haar-bessel}, and the operator bound,
\[
 \sum_{Q\in\mathcal A}\mu(Q)
 \le H_I I C_{\rm R}^2\mu(S)
 \le H_I I C_{\rm R}^2
       \frac{G(C_0+2)^n}{c_1}\mu(Q_0).
\]
The empty family is immediate. Taking the supremum over finite
subfamilies proves that $\mathcal F_1$ is Carleson. Finally add the possible top level with $\ell(Q)>D_E$ to obtain
$\mathcal F$. Its mass below any top cube is at most the mass of that
top cube. The seed conclusion and its uniform depth bound
hold for every cube outside this family.
\end{proof}
